\documentclass[preprint,12pt]{elsarticle}
\journal{Knowledge-Based Systems}
\usepackage{graphicx,booktabs,amsmath,amssymb}
\newcommand{\E}{\mathbb{E}}
\newcommand{\candidates}{\mathcal{C}}
\newcommand{\loss}{\ell}
\newcommand{\marginal}{m}
\usepackage{hyperref}
\setlength{\emergencystretch}{1em}
\begin{document}
\begin{frontmatter}
\title{Supplementary Material: Response-Aware Marginal-Utility Routing for Multi-Context Prediction}
\ifdefined\anonymousmode
\author{}
\else
\author[zua]{Xinling Wen}
\author[zua]{Jiahao Zhang}
\author[zua]{Yu Chen}
\affiliation[zua]{organization={School of Electronics and Information,
                             Zhengzhou University of Aeronautics},
                 city={Zhengzhou}, country={China}}
\fi
\begin{abstract}
This supplement documents the controlled environment, the implementation and
training configuration of R-MUR, the multi-target traffic protocol, the
complete per-seed results, and the additional response-observability analysis
that supports the mechanism section of the main paper.
\end{abstract}
\end{frontmatter}

\section{Controlled Environment}

\subsection{Data generator}
Each episode contains a target vector $y\in\mathbb{R}^{G}$, an anchor prediction
$a=y+\sigma_a\varepsilon$ with $\varepsilon\sim\mathcal{N}(0,I)$, and $K$
candidate contexts. The candidate pool is built from
$U=\max(1,\operatorname{round}(K(1-r_{\mathrm{dup}})))$ distinct groups, where
$r_{\mathrm{dup}}$ is the duplicate ratio. Each candidate is assigned to one of
these groups; duplicates share the same group assignment. For a candidate in
group $g$, the useful measurement is $m_g=y_g+\rho_{\mathrm{use}}(a_g-y_g)$ with
$\rho_{\mathrm{use}}=0.1$. A candidate marked harmful receives
$m_g=y_g+\rho_{\mathrm{harm}}(a_g-y_g)$ with $\rho_{\mathrm{harm}}=1.8$. The
frozen expert predicts each group by the mean measurement of its selected
candidates, leaving groups without a selected candidate at their anchor value.
The episode loss is the mean squared error across groups. Standalone and
state-conditioned marginals are computed by enumerating every candidate
addition. Relevance is drawn independently from $[0.7,1.0]$; duplicate groups
therefore make content redundancy observable to similarity-based selectors.
Maximum marginal relevance follows the standard greedy rule with relevance
minus a cosine-similarity penalty to the selected set; the penalty is selected
on a validation split from $\{0,.25,.5,1\}$.

\subsection{Protocols}
\begin{table}[t]
\centering
\small
\caption{Controlled protocols. All reported gains are loss reductions relative
to the anchor-only predictor. Seeds are independent training and evaluation
runs.}
\resizebox{\linewidth}{!}{%
\begin{tabular}{lllll}
\toprule
Protocol & $K$ & $B$ & $r_{\mathrm{dup}}$ & Seeds / splits\\
\midrule
Redundancy ladder & 16 & 4 & $0,.25,.5,.75$ & 5 seeds; 3000 train, 2500 test\\
Candidate scaling & $4,8,16,32,64$ & 4 & .50 & 5 seeds; 2000 train, 1500 test\\
\bottomrule
\end{tabular}}
\end{table}

The controlled experiments use 15 training epochs for the candidate-scaling
protocol and 20 epochs for the redundancy ladder.

\section{R-MUR Implementation and Training}
\label{sec:supp-rmur}

\subsection{Router}
The router has three parts: a cached scorer $g_C$, a response-aware scorer
$g_R$, and one shared set encoder. Every candidate is represented by its
history, the query by the anchor history, and the selected set by the
histories of the selected candidates packed into a fixed-size tensor with a
validity mask. A permutation-invariant DeepSets encoder of depth two with
hidden and set width 64 summarises the selected set. The cached scorer maps the
query, set and candidate representations to a scalar. The response-aware
scorer takes the same inputs plus the seven response features described below
and is a two-layer network of width 64.

\subsection{Response summary}
For a candidate $j$ at state $A$ the frozen predictor is evaluated with and
without the candidate, and the summary concatenates: the mean of the base
forecast, the mean of the forecast after adding the candidate, the mean of the
change, its standard deviation, its mean absolute value, its Euclidean norm
divided by $\sqrt{H}$ with horizon $H$, and its maximum absolute value.

\subsection{Training}
Training states are sampled from the training episodes. A state is a uniformly
random subset of the candidate pool whose size is drawn uniformly from $0$ to
$B-1$, so empty, partial and nearly complete selected sets all occur. Every
candidate outside the state contributes one labelled example with its true
marginal utility, computed as the counterfactual loss difference between the
predictor with and without that candidate. A single state therefore supplies
$K-|A|$ regression targets and all candidate pairs of the state for the ranking
term. The objective is a Huber regression term plus a pairwise ranking term
whose per-pair weight is the absolute utility gap of the pair,
\begin{equation}
 \mathcal L
 =
 \mathcal L_{\mathrm{Huber}}
 +
 \beta
 \sum_{i,j:\,\widetilde m_i>\widetilde m_j}
 |\widetilde m_i-\widetilde m_j|
 \log\!\left(1+\exp\!\left[-(\widehat m_i-\widehat m_j)\right]\right),
\end{equation}
with $\beta=1$ in every reported run. Training uses AdamW with learning rate
$10^{-3}$ and weight decay $10^{-4}$, batches of 128 states, and 20 epochs;
each target-run contributes 16{,}000 states, four per training episode, with 16
candidate labels each.

\subsection{Inference}
At deployment the cached scorer ranks the remaining candidates and the
response-aware scorer is evaluated only on the top-$q$ shortlist, with $q=4$ in
every reported run. The router adds the highest-scoring shortlisted candidate
when its estimated marginal utility is positive and otherwise stops; it also
stops when the context budget is reached. The predictor is therefore evaluated
on a fraction $q/K$ of the pool per decision.

\subsection{Baselines}
\label{sec:supp-baselines}
Standalone Utility is a model with the same set encoder, trained on states with
an empty selected set and one sampled candidate per state, so it estimates the
standalone marginal $u(j\mid\emptyset)$ and its ranking is fixed before
selection. The cached-only router is the same model trained on random nonempty
states from cached inputs only, so it is state-conditioned but never observes
the predictor response. Both use the same optimiser, schedule and state count
as R-MUR.

\paragraph{Content-based subset rules.}
Facility location maximises the sum of similarities to the target, greedy DPP
maximises the determinant of the candidate kernel, facility location with
relevance weights multiplies the two, and $k$-center minimises the largest
distance from the target to the selected set. All four use the same candidate
pool and budget and require no training.

\paragraph{Fixed-set acquisition rules.}
Five policies from the feature-selection and active-feature-acquisition
literature are adapted to the protocol. Every policy uses only the training and
validation splits: the harm rate is the per-episode fraction of episodes whose
loss exceeds the anchor-only loss, the same definition used for every other
policy in the paper. Each chooses one set per target from
the training and validation splits with the frozen predictor, and then applies
that set to every test episode; none of them observes a test label, and none
conditions its decision on the state of an individual episode. They answer a
question the router must answer, namely whether per-episode conditioning buys
anything over the best set that can be chosen offline.

\begin{itemize}
\item \emph{Greedy validation set} follows best-subset and optimal-pursuit
  selection \citep{zhu2025best}: starting from the empty set, it adds the
  candidate with the largest validation loss reduction, and one elimination
  pass then removes the weakest element and reconsiders its best replacement.
  The elimination pass selected the same set as plain forward selection on all
  576 target-runs of the sweep, so the two are not reported separately.
\item \emph{Adaptive-$k$ relevance} follows adaptive-$k$ context selection
  \citep{taguchi2025adaptive}: candidates are ordered by relevance and the
  number of contexts is chosen on validation rather than fixed in advance, so
  the policy may select fewer than the budget when additional contexts do not
  reduce the validation loss.
\item \emph{Cluster-based selection} follows iterative demonstration
  selection \citep{qin2023iterative}: the candidate pool is clustered in
  history space and the most relevant member of each cluster represents it, so
  the chosen set is diverse by construction and needs no predictor. The cluster
  structure is estimated on the validation split; no statistic of the test
  episodes enters the selection, and the fixed set is then applied to them.
\item \emph{Two-step lookahead set} follows nongreedy active feature
  acquisition \citep{valancius2023acquisition,norcliffe2025stochastic}: the
  best pair of candidates is chosen on validation before any single candidate
  is committed to, and the set is then completed greedily, so jointly
  informative candidates are not missed by a myopic rule. The policy selects
  the same set as greedy validation selection on most runs and matches its gain
  to within .0002 on every dataset, so non-greedy lookahead does not change the
  outcome in this setting.
\item \emph{Gradient influence} ranks candidates by the first-order loss
  reduction, so the sign convention matters: $2\langle d_j, r\rangle+\lVert
  d_j\rVert^2$ is the increase in loss and the ranking uses its negation. The
  policy is scored on the validation split and the resulting set is applied to
  the test episodes.
\item \emph{Gradient influence} follows gradient-estimated demonstration
  selection \citep{zhang2025linear}: the first-order term of the squared-loss
  change, $2\langle f_A-y, d_j\rangle+\lVert d_j\rVert^2$ with
  $d_j=f_{A\cup\{j\}}-f_A$, scores each candidate on the validation split,
  averaged over randomly sampled states so that the score does not depend on
  one state; candidates are scored independently, in linear time in the pool
  size. Because that expression is a loss increase, the ranking selects the
  candidates with the smallest value.
\end{itemize}

The oracle policies use the true marginals and are not deployable.

\section{Extensions of the Regret Decomposition}
\label{sec:supp-proofs}

The main paper states the shortlist-aware bound
\begin{equation}
 \marginal(j^\star\mid A)-\marginal(\widehat{j}\mid A)\le 2\epsilon_C+2\epsilon_R ,
\end{equation}
which follows from the exact decomposition into a screening and a reranking
term. One remark extends it.

\paragraph{Sequential selection.}
A guarantee for a sequence of $B$ selections requires the induced set utility
to be monotone and submodular, so that approximate greedy selection inherits
the standard $(1-1/e)$ factor up to the accumulated estimation error. Context
utility is not monotone in general, because a context can increase the loss,
and we do not assume submodularity in the experiments. The sequential case is
therefore stated as a condition rather than claimed for R-MUR.

\section{Multi-Target Traffic Protocol}

The six benchmarks use the standard files distributed with the traffic
forecasting benchmark suite. Each file provides a multivariate series and
fixed train, validation, and test index sets; no observation is reordered
across the split boundaries. Values are standardized with the mean and standard
deviation of the training segment.

For a dataset with $N$ sensors, the evaluation fixes 32 evenly spaced target
sensors. Each target sensor has its own candidate pool, built from the training
segment: eight sensors with the highest correlation to the target, five sensors
of intermediate correlation, and three sensors drawn uniformly at random with a
fixed construction seed. Each episode takes a window from the split under
evaluation. The anchor is the target sensor's previous 12 observations, each
candidate contributes its previous 12 observations, and the target is the
target sensor's next 12 observations.

The frozen predictor is a ridge model with penalty 10, fitted for each target
sensor on that target's training windows under the full context subset, the
empty subset, and a randomly masked subset of its own candidate pool, and then
held fixed for candidate evaluation, oracle computation and router training.
The predictor is therefore specific to the target and is refitted for every
target. The loss is the mean squared error over the forecast horizon.

Each target is repeated with three independent runs. The reported prediction
gain is the reduction in loss relative to the anchor-only predictor, the
negative-transfer rate is the fraction of episodes whose routed loss exceeds
the anchor-only loss, and oracle recovery is the pooled ratio of learned gain
to sequential-oracle gain. Intervals resample the 32 target means with 5000
bootstrap replicates; the hierarchical variant resamples targets and then runs
within each target.

\section{Complete Controlled and Multi-Target Results}

\begin{table}[t]
\centering
\small
\caption{Per-seed controlled redundancy ladder at $K=16$ and $B=4$. Gains are loss reductions relative to the anchor-only predictor.}
\label{tab:supp-redundancy}
\resizebox{\linewidth}{!}{%
\begin{tabular}{lccccc}
\toprule
$r_{\mathrm{dup}}$ & Seed & Oracle-Static & R-MUR & Oracle-Greedy & $H_{\mathrm{state}}$\\
\midrule
0 & 101 & .691 & .663 & .691 & .000 \\
0 & 202 & .678 & .655 & .678 & .000 \\
0 & 303 & .697 & .677 & .697 & .000 \\
0 & 404 & .692 & .669 & .692 & .000 \\
0 & 505 & .681 & .653 & .681 & .000 \\
0.25 & 101 & .517 & .502 & .585 & .115 \\
0.25 & 202 & .509 & .489 & .574 & .113 \\
0.25 & 303 & .518 & .510 & .584 & .114 \\
0.25 & 404 & .523 & .516 & .588 & .111 \\
0.25 & 505 & .513 & .497 & .575 & .108 \\
0.5 & 101 & .353 & .369 & .450 & .216 \\
0.5 & 202 & .347 & .361 & .442 & .215 \\
0.5 & 303 & .353 & .366 & .449 & .213 \\
0.5 & 404 & .350 & .366 & .451 & .224 \\
0.5 & 505 & .353 & .364 & .444 & .204 \\
0.75 & 101 & .176 & .206 & .247 & .286 \\
0.75 & 202 & .175 & .205 & .244 & .283 \\
0.75 & 303 & .179 & .208 & .250 & .282 \\
0.75 & 404 & .179 & .217 & .253 & .293 \\
0.75 & 505 & .176 & .207 & .246 & .285 \\
\bottomrule
\end{tabular}}
\end{table}

\begin{table}[t]
\centering
\small
\caption{Per-seed candidate-count scaling under fixed redundancy. Normalized MAE divides pointwise error by the standard deviation of the true marginal targets. Decision accuracy counts a selection as correct when it attains the best true utility, including exact ties.}
\label{tab:supp-candidate-scaling}
\resizebox{\linewidth}{!}{%
\begin{tabular}{rcccccccc}
\toprule
$K$ & Seed & R-MUR recovery & Normalized MAE & Positive margin & Strict top-1 & Decision accuracy & One-step regret\\
\midrule
4 & 101 & .961 & .387 & .205 & .470 & .873 & .0192 \\
4 & 202 & .960 & .412 & .230 & .493 & .887 & .0164 \\
4 & 303 & .969 & .426 & .216 & .491 & .890 & .0167 \\
4 & 404 & .979 & .382 & .230 & .495 & .887 & .0201 \\
4 & 505 & .969 & .408 & .221 & .492 & .883 & .0186 \\
8 & 101 & .799 & .612 & .116 & .227 & .383 & .0654 \\
8 & 202 & .809 & .565 & .111 & .237 & .403 & .0573 \\
8 & 303 & .800 & .613 & .113 & .228 & .378 & .0666 \\
8 & 404 & .843 & .556 & .112 & .252 & .416 & .0535 \\
8 & 505 & .815 & .574 & .116 & .240 & .396 & .0608 \\
16 & 101 & .687 & .698 & .071 & .174 & .285 & .0693 \\
16 & 202 & .699 & .681 & .066 & .183 & .298 & .0656 \\
16 & 303 & .709 & .687 & .071 & .185 & .301 & .0650 \\
16 & 404 & .708 & .679 & .072 & .182 & .299 & .0667 \\
16 & 505 & .684 & .684 & .071 & .176 & .289 & .0687 \\
32 & 101 & .604 & .751 & .042 & .141 & .224 & .0608 \\
32 & 202 & .600 & .733 & .041 & .141 & .220 & .0588 \\
32 & 303 & .591 & .718 & .044 & .140 & .222 & .0618 \\
32 & 404 & .586 & .728 & .044 & .137 & .222 & .0626 \\
32 & 505 & .598 & .728 & .040 & .140 & .219 & .0586 \\
64 & 101 & .470 & .718 & .025 & .089 & .136 & .0513 \\
64 & 202 & .484 & .731 & .025 & .098 & .151 & .0482 \\
64 & 303 & .479 & .756 & .023 & .096 & .155 & .0479 \\
64 & 404 & .486 & .739 & .022 & .088 & .149 & .0484 \\
64 & 505 & .480 & .751 & .023 & .094 & .146 & .0483 \\
\bottomrule
\end{tabular}}
\end{table}

\begin{table}[t]
\centering
\small
\caption{Per-seed METR-LA observability probe. $\mathcal{X}_1$ uses cached context features, $\mathcal{X}_2$ adds scalar response summaries, and $\mathcal{X}_3$ adds full forecast trajectories. MAE is the mean absolute error and AUC is the area under the receiver operating characteristic curve.}
\label{tab:supp-observability}
\resizebox{\linewidth}{!}{%
\begin{tabular}{llccccc}
\toprule
Seed & Input & MAE & Spearman & Positive AUC & Top-1 & Regret\\
\midrule
101 & $\mathcal{X}_1$ & .0339 & .169 & .651 & .173 & .0635 \\
101 & $\mathcal{X}_2$ & .0308 & .377 & .829 & .260 & .0503 \\
101 & $\mathcal{X}_3$ & .0307 & .333 & .839 & .242 & .0505 \\
202 & $\mathcal{X}_1$ & .0400 & .162 & .649 & .158 & .0677 \\
202 & $\mathcal{X}_2$ & .0351 & .383 & .838 & .208 & .0538 \\
202 & $\mathcal{X}_3$ & .0353 & .383 & .843 & .237 & .0531 \\
303 & $\mathcal{X}_1$ & .0418 & .144 & .655 & .124 & .0796 \\
303 & $\mathcal{X}_2$ & .0370 & .391 & .848 & .225 & .0580 \\
303 & $\mathcal{X}_3$ & .0376 & .363 & .855 & .223 & .0623 \\
404 & $\mathcal{X}_1$ & .0405 & .172 & .640 & .131 & .0807 \\
404 & $\mathcal{X}_2$ & .0371 & .390 & .828 & .285 & .0667 \\
404 & $\mathcal{X}_3$ & .0370 & .391 & .835 & .262 & .0654 \\
505 & $\mathcal{X}_1$ & .0427 & .171 & .652 & .154 & .0765 \\
505 & $\mathcal{X}_2$ & .0384 & .411 & .827 & .320 & .0559 \\
505 & $\mathcal{X}_3$ & .0384 & .393 & .834 & .258 & .0563 \\
\bottomrule
\end{tabular}}
\end{table}

\begin{table}[t]
\centering
\small
\caption{Structural headroom per dataset over 32 fixed targets and five runs (960 target-run evaluations). Every evaluation has positive headroom.}
\label{tab:supp-headroom}
\resizebox{\linewidth}{!}{%
\begin{tabular}{lrrr}
\toprule
Dataset & Mean headroom & Median headroom & Positive fraction\\
\midrule
METRLA & .097 & .088 & 1.000 \\
PEMSBAY & .115 & .116 & 1.000 \\
PEMS03 & .122 & .111 & 1.000 \\
PEMS04 & .085 & .076 & 1.000 \\
PEMS07 & .089 & .087 & 1.000 \\
PEMS08 & .078 & .077 & 1.000 \\
\bottomrule
\end{tabular}}
\end{table}

\begin{table}[t]
\centering
\small
\caption{Per-seed multi-target prediction gains for the learned policies, averaged over the 32 fixed targets of each dataset.}
\label{tab:supp-multitarget}
\resizebox{\linewidth}{!}{%
\begin{tabular}{llrrr}
\toprule
Dataset & Run & Standalone & Cached-only & R-MUR\\
\midrule
METRLA & 101 & .001 & .002 & .015 \\
METRLA & 202 & .001 & .006 & .012 \\
METRLA & 303 & .005 & .005 & .017 \\
PEMSBAY & 101 & .023 & .027 & .055 \\
PEMSBAY & 202 & .027 & .028 & .052 \\
PEMSBAY & 303 & .031 & .030 & .060 \\
PEMS03 & 101 & .015 & .015 & .027 \\
PEMS03 & 202 & .016 & .017 & .029 \\
PEMS03 & 303 & .014 & .014 & .025 \\
PEMS04 & 101 & .024 & .025 & .035 \\
PEMS04 & 202 & .025 & .024 & .035 \\
PEMS04 & 303 & .025 & .025 & .037 \\
PEMS07 & 101 & .014 & .014 & .028 \\
PEMS07 & 202 & .017 & .016 & .029 \\
PEMS07 & 303 & .014 & .015 & .028 \\
PEMS08 & 101 & .025 & .025 & .034 \\
PEMS08 & 202 & .024 & .024 & .033 \\
PEMS08 & 303 & .024 & .025 & .033 \\
\bottomrule
\end{tabular}}
\end{table}

\begin{table}[t]
\centering
\small
\caption{Content-based subset-selection baselines under the multi-target protocol, averaged over 32 fixed targets and three runs. None of these rules trains a model or evaluates the predictor during selection.}
\label{tab:supp-subset-baselines}
\resizebox{\linewidth}{!}{%
\begin{tabular}{lrrrr}
\toprule
Dataset & Facility location & DPP greedy & Facility loc. + rel. & $k$-center\\
\midrule
METRLA & -.021 & -.015 & -.017 & -.018 \\
PEMSBAY & .008 & .005 & -.007 & -.014 \\
PEMS03 & .009 & .007 & -.002 & .004 \\
PEMS04 & .014 & .013 & .012 & .012 \\
PEMS07 & .007 & .006 & -.002 & .003 \\
PEMS08 & .017 & .012 & .013 & .009 \\
\midrule
Mean & .006 & .005 & -.000 & -.000 \\
\bottomrule
\end{tabular}}
\end{table}


\section{Additional Response-Observability Analysis}

The probe behind the mechanism panel of the main paper uses the frozen ridge
predictor of a single target sensor of METR-LA and five independent runs. It
samples four states per episode. For each state and candidate the feature
families are:
\begin{itemize}
\item $\mathcal{X}_1$: query history, mean selected history, candidate history,
their difference and product, and selected-set size.
\item $\mathcal{X}_2$: $\mathcal{X}_1$ plus the mean expert prediction before
and after adding the candidate, the mean prediction delta, and the mean
absolute delta.
\item $\mathcal{X}_3$: $\mathcal{X}_2$ plus the full before and after forecast
trajectories and their difference.
\end{itemize}
A histogram-based gradient boosting regressor and classifier are fitted for
each family, with hyperparameters selected on a calibration split over depth
$\{3,6,\infty\}$ and $L_2$ regularization $\{10^{-3},1\}$, 150 boosting
iterations and learning rate .08. The classifier provides the positive-utility
AUC; the regressor provides the mean absolute error, Spearman correlation,
top-1 accuracy and one-step regret reported in the main paper. The probe is a
diagnostic of the information content of the predictor response, not a routing
policy.

\section{Decision-Margin Audit}

The candidate-scaling audit uses the saved rollout predictions and true
marginals of the controlled scaling protocol. For each state and episode, let
$m_{(1)}$ and $m_{(2)}$ be the largest and second-largest true marginal
utilities among the unselected candidates. The positive top-two margin is
$\Delta=m_{(1)}-m_{(2)}$ when $\Delta>0$. We report the normalized pointwise
error
\begin{equation}
 \mathrm{nMAE}=\frac{\mathbb{E}|\widehat m-m|}{\operatorname{SD}(m)},
\end{equation}
the fraction of states with zero top-two margin, the median ratio between the
maximum absolute candidate error and the positive margin, and the decision
accuracy
\begin{equation}
 \mathrm{Acc}_{\mathrm{dec}}=
 \Pr\!\left[m(\widehat j\mid A)=m(j^*\mid A)\right],
\end{equation}
where $\widehat j$ is the predicted argmax and $j^*$ is a true argmax. Strict
top-1 accuracy compares candidate indices, so it counts an exact tie among
equally useful candidates as an error. Accuracy and margin statistics use
states with at least two unselected candidates, where a nontrivial argmax
decision exists.

\end{document}
